\section{Discussion and Limitations}

\paragraph{Relationship to KVFlow and KVCOMM.}
\system uses KVFlow and \kvcomm as reference designs, not as claimed contributions.
The contribution is the coding-specific layer around them: workflow-derived code segment hints, exact content signatures, conservative safety gating, and SWE-bench-oriented validation.
This distinction matters because the same underlying cache scheduler or cross-context reuse primitive could be replaced by another serving backend without changing the code-base safety invariant.

\paragraph{Lossy does not mean approximate code reuse.}
Reuse is lossy only because the reused KV tensors are transformed across positions.
The code content must be exactly identical.
This distinction is central to the safety argument and should be preserved in any deployment.
In the 100-case serving scalability experiment (Table~\ref{tab:prefetch}), exact-content reuse achieves a 99\% hit rate with 64\% more cached tokens than the prefix-only baseline, confirming that lossy-position reuse does not compromise code content integrity.

\paragraph{Patch quality is a separate bottleneck.}
The 30-case pass@1 run uses a 7B instruction model and constrained JSON edits.
The absolute pass@1 is low for both lossless and lossy modes.
We therefore interpret the accuracy result as a lossless-vs-lossy delta test, not as a claim that the current patch generator is state of the art.
A stronger coder model or richer repair loop should improve absolute pass@1 while preserving the same KV reuse comparison.

\paragraph{Hardware-constrained model evaluation.}
We attempted to evaluate a 32B-parameter quantized model (Qwen2.5-32B-Instruct-GPTQ-Int4) on a single RTX~4090 (24\,GB) as paired accuracy evidence.
The model weights consume nearly all GPU memory, leaving only ${\sim}2{,}554$ input-token capacity---insufficient for realistic repo-level coding prompts.
Auto-truncation causes an internal metadata index error in the serving runtime.
We therefore report all accuracy and serving results with the 7B model and acknowledge that stronger-model validation requires multi-GPU or higher-memory hardware.

\paragraph{Host-backed prefetch needs more engineering.}
The file-backed HiCache storage backend triggers a token-to-KV-pool allocator leak during serving, causing runtime crashes.
The present coding-aware prefetch result (Table~\ref{tab:prefetch}) is therefore obtained with host-backed storage disabled, validating the interface, metadata path, and exact-content reuse behavior but not host load-back acceleration.
Completing robust host load-back is future work for stronger TTFT gains.

\paragraph{Scope.}
\system targets coding MAS prompts with large repeated code-base segments.
It is not designed to reuse arbitrary natural-language context, approximate code variants, or semantically equivalent but textually different implementations.
